% ===========================================================================
\chapter{Especificação dos Casos de Uso de Segurança e Parametrização Clínica}
\label{cap_seguranca_parametrizacao}
% ===========================================================================

Este capítulo especifica formalmente os casos de uso compreendidos nos pacotes funcionais de Segurança e Controle de Acesso e de Parametrização Clínica da Metodologia IHI-GTT. A Figura \ref{fig_uc_modulo_seguranca} ilustra o diagrama de casos de uso deste subsistema em notação UML com renderização em TikZ:

\begin{figure}[!ht]
\centering
\caption{Diagrama de Casos de Uso do Subsistema de Segurança e Parametrização}
\label{fig_uc_modulo_seguranca}
\resizebox{0.90\textwidth}{!}{%
\begin{tikzpicture}[
    actor/.style={
        draw=clinical-900,
        fill=clinical-100!40,
        thick,
        rounded corners=2pt,
        align=center,
        font=\bfseries\scriptsize,
        minimum width=2.4cm,
        minimum height=1.2cm
    },
    usecase/.style={
        ellipse,
        draw=clinical-700,
        fill=white,
        thick,
        align=center,
        font=\scriptsize,
        minimum width=3.4cm,
        minimum height=1.1cm,
        inner sep=2pt
    },
    link/.style={
        draw=clinical-900!80,
        thick
    },
    extend_link/.style={
        draw=alert-amber!90!black,
        thick,
        dashed,
        -Latex
    }
]
    % Fronteira
    \draw[rounded corners=8pt, fill=clinical-100!15, draw=clinical-700, line width=1.2pt] 
        (0.5, -7.2) rectangle (12.5, 3.2);
    \node[anchor=north west, font=\bfseries\small\color{clinical-700}] at (0.8, 2.9) 
        {Subsistema: Segurança e Parametrização GTT};

    % Atores
    \node[actor] (admin) at (-2.2, 1.5) {Administrador\\\texttt{\scriptsize (ADMIN)}};
    \node[actor] (prof) at (-2.2, -2.0) {Docente\\\texttt{\scriptsize (PROFESSOR)}};
    \node[actor] (student) at (-2.2, -5.5) {Discente\\\texttt{\scriptsize (STUDENT)}};

    % Casos de Uso
    \node[usecase] (uc01) at (3.5, 1.5) {UC01: Autenticar Usuário\\no Sistema};
    \node[usecase] (uc02) at (8.5, 1.5) {UC02: Forçar Redefinição\\de Senha Provisória};
    \node[usecase] (uc03) at (3.5, -0.5) {UC03: Alterar Senha\\pelo Perfil};
    \node[usecase] (uc04) at (8.5, -0.5) {UC04: Gerenciar Contas\\de Usuários};
    \node[usecase] (uc05) at (3.5, -3.5) {UC05: Parametrizar Módulos\\e Gatilhos GTT};
    \node[usecase] (uc06) at (8.5, -3.5) {UC06: Consultar Catálogo\\GTT e Gravidades};

    % Conexões
    \draw[link] (admin) -- (uc01);
    \draw[link] (admin) -- (uc03);
    \draw[link] (admin) -- (uc04);
    \draw[link] (admin) -- (uc05);
    \draw[link] (admin) -- (uc06);

    \draw[link] (prof) -- (uc01);
    \draw[link] (prof) -- (uc03);
    \draw[link] (prof) -- (uc06);

    \draw[link] (student) -- (uc01);
    \draw[link] (student) -- (uc03);
    \draw[link] (student) -- (uc06);

    % Extensao
    \draw[extend_link] (uc02) -- node[above, font=\tiny\bfseries] {<<extend>>} (uc01);

\end{tikzpicture}
}
\fonte{Elaborado pelo autor (2026).}
\end{figure}

% ===========================================================================
\section{Especificação Detalhada dos Casos de Uso}
\label{sec_especificacao_uc_seguranca}
% ===========================================================================

% ---------------------------------------------------------------------------
% UC01
% ---------------------------------------------------------------------------
\subsection{UC01: Autenticar Usuário no Sistema}
\label{uc01_especificacao}

\begin{table}[!ht]
\centering
\caption{Especificação Detalhada do Caso de Uso UC01}
\label{tab_uc01}
\footnotesize
\begin{tabularx}{\textwidth}{p{3.5cm} X}
\toprule
\textbf{Propriedade} & \textbf{Especificação Técnica e Comportamental} \\
\midrule
\textbf{Identificador e Nome} & \texttt{UC01: Autenticar Usuário no Sistema} \\
\textbf{Atores} & \texttt{ADMIN}, \texttt{PROFESSOR}, \texttt{STUDENT} (Atores Primários). \\
\textbf{Nível de Meta} & Nível Usuário (\textit{Sea Level}). \\
\textbf{Pré-condições} & O usuário deve possuir cadastro prévio ativo no banco de dados. \\
\textbf{Gatilho} & O usuário acessa a URL raiz da plataforma e insere suas credenciais. \\
\textbf{Fluxo Principal} & 
1. O usuário informa seu e-mail institucional e sua senha.\\
& 2. O sistema valida se o e-mail pertence estritamente ao domínio \texttt{@academico.ufs.br}.\\
& 3. O sistema recupera o registro na tabela \texttt{users} e verifica a senha com BCrypt.\\
& 4. O sistema verifica que o usuário está ativo (\texttt{is\_active = true}).\\
& 5. O sistema verifica que a flag \texttt{must\_change\_password} é \texttt{false}.\\
& 6. O sistema emite um token JWT \textit{stateless} contendo as reivindicações do perfil.\\
& 7. O sistema direciona o usuário para o painel de bordo (\textit{dashboard}) de seu perfil. \\
\textbf{Fluxos Alternativos} & 
\textbf{[5a] Senha Provisória Pendente}: O sistema identifica que \texttt{must\_change\_password = true}. O sistema aciona o ponto de extensão e invoca \texttt{UC02: Forçar Redefinição de Senha Provisória}. \\
\textbf{Fluxos de Exceção} & 
\textbf{[2e] E-mail Não Acadêmico}: Se o domínio for diferente de \texttt{@academico.ufs.br}, o sistema rejeita a tentativa antes de consultar o banco, exibindo mensagem de erro amigável.\\
& \textbf{[3e] Credenciais Inválidas}: Se a senha não conferir com o hash armazenado, o sistema emite mensagem genérica de falha de autenticação e bloqueia o acesso.\\
& \textbf{[4e] Conta Inativa}: Se \texttt{is\_active = false}, o sistema informa que a conta está suspensa. \\
\textbf{Pós-condições} & 
\textit{Sucesso}: Token JWT gerado e sessão ativa na estação de trabalho.\\
& \textit{Falha}: Nenhuma credencial é emitida e o acesso permanece bloqueado. \\
\textbf{Regras de Negócio} & \textbf{RN01} (Domínio Exclusivo Institucional) e \textbf{RN04} (Criptografia BCrypt). \\
\textbf{Requisitos Associados} & \texttt{RF01}, \texttt{RF04}, \texttt{RNF01}, \texttt{RNF02}. \\
\bottomrule
\end{tabularx}
\fonte{Elaborado pelo autor (2026).}
\end{table}

% ---------------------------------------------------------------------------
% UC02
% ---------------------------------------------------------------------------
\subsection{UC02: Forçar Redefinição de Senha Provisória}
\label{uc02_especificacao}

\begin{table}[!ht]
\centering
\caption{Especificação Detalhada do Caso de Uso UC02}
\label{tab_uc02}
\footnotesize
\begin{tabularx}{\textwidth}{p{3.5cm} X}
\toprule
\textbf{Propriedade} & \textbf{Especificação Técnica e Comportamental} \\
\midrule
\textbf{Identificador e Nome} & \texttt{UC02: Forçar Redefinição de Senha Provisória} \\
\textbf{Atores} & \texttt{ADMIN}, \texttt{PROFESSOR}, \texttt{STUDENT} (Atores Primários). \\
\textbf{Nível de Meta} & Nível Subfunção / Extensão (\textit{Fish Level}). \\
\textbf{Pré-condições} & O usuário foi autenticado com sucesso em \texttt{UC01} com \texttt{must\_change\_password = true}. \\
\textbf{Gatilho} & Disparo pelo ponto de extensão de \texttt{UC01} após validação de credenciais provisórias. \\
\textbf{Fluxo Principal} & 
1. O sistema redireciona a interface forçadamente para a tela de troca de senha.\\
& 2. O usuário informa sua nova senha e a confirmação de senha.\\
& 3. O sistema valida os critérios de complexidade da senha (mínimo de 8 caracteres).\\
& 4. O sistema gera o novo hash BCrypt e atualiza a tabela \texttt{users}.\\
& 5. O sistema atualiza a flag \texttt{must\_change\_password} para \texttt{false}.\\
& 6. O sistema emite notificação \textit{Toast} de sucesso e redireciona ao painel de bordo. \\
\textbf{Fluxos de Exceção} & 
\textbf{[3e] Critérios de Complexidade Não Atendidos}: Se a senha possuir menos de 8 caracteres ou confirmação divergente, o sistema rejeita o formulário e exige correção.\\
& \textbf{[4e] Tentativa de Burla de Navegação}: Se o usuário tentar acessar qualquer outra rota via barra de endereços, o \textit{AuthGuard} do frontend intercepta a rota e força o retorno à tela de troca de senha. \\
\textbf{Pós-condições} & Senha atualizada e flag provisória cancelada. Acesso regular liberado. \\
\textbf{Regras de Negócio} & \textbf{RN02} (Obrigatoriedade de Troca de Senha no Primeiro Login). \\
\textbf{Requisitos Associados} & \texttt{RF02}, \texttt{RNF01}. \\
\bottomrule
\end{tabularx}
\fonte{Elaborado pelo autor (2026).}
\end{table}

% ---------------------------------------------------------------------------
% UC03
% ---------------------------------------------------------------------------
\subsection{UC03: Alterar Senha de Acesso pelo Perfil}
\label{uc03_especificacao}

\begin{table}[!ht]
\centering
\caption{Especificação Detalhada do Caso de Uso UC03}
\label{tab_uc03}
\footnotesize
\begin{tabularx}{\textwidth}{p{3.5cm} X}
\toprule
\textbf{Propriedade} & \textbf{Especificação Técnica e Comportamental} \\
\midrule
\textbf{Identificador e Nome} & \texttt{UC03: Alterar Senha de Acesso pelo Perfil} \\
\textbf{Atores} & Qualquer usuário autenticado (\texttt{ADMIN}, \texttt{PROFESSOR}, \texttt{STUDENT}). \\
\textbf{Nível de Meta} & Nível Usuário (\textit{Sea Level}). \\
\textbf{Pré-condições} & O usuário está autenticado no sistema com token JWT válido. \\
\textbf{Gatilho} & O usuário acessa as configurações de seu perfil e clica em "Alterar Senha". \\
\textbf{Fluxo Principal} & 
1. O usuário informa a senha atual, a nova senha e a confirmação.\\
& 2. O sistema verifica a exatidão da senha atual contra o hash do banco de dados.\\
& 3. O sistema valida as regras de complexidade da nova senha.\\
& 4. O sistema armazena o novo hash BCrypt e registra carimbo \texttt{updated\_at}.\\
& 5. O sistema exibe notificação de confirmação e mantém a sessão ativa. \\
\textbf{Fluxos de Exceção} & 
\textbf{[2e] Senha Atual Incorreta}: O sistema informa erro de verificação e interrompe a troca.\\
& \textbf{[3e] Nova Senha Idêntica à Atual}: O sistema exige que a nova credencial seja distinta. \\
\textbf{Pós-condições} & Senha alterada com sucesso no banco de dados. \\
\textbf{Regras de Negócio} & \textbf{RN04} (Criptografia de Credenciais). \\
\textbf{Requisitos Associados} & \texttt{RF03}, \texttt{RNF01}. \\
\bottomrule
\end{tabularx}
\fonte{Elaborado pelo autor (2026).}
\end{table}

% ---------------------------------------------------------------------------
% UC04
% ---------------------------------------------------------------------------
\subsection{UC04: Gerenciar Contas de Usuários}
\label{uc04_especificacao}

\begin{table}[!ht]
\centering
\caption{Especificação Detalhada do Caso de Uso UC04}
\label{tab_uc04}
\footnotesize
\begin{tabularx}{\textwidth}{p{3.5cm} X}
\toprule
\textbf{Propriedade} & \textbf{Especificação Técnica e Comportamental} \\
\midrule
\textbf{Identificador e Nome} & \texttt{UC04: Gerenciar Contas de Usuários} \\
\textbf{Atores} & \texttt{ADMIN} (Ator Primário). \\
\textbf{Nível de Meta} & Nível Usuário (\textit{Sea Level}). \\
\textbf{Pré-condições} & O administrador está autenticado com perfil \texttt{ADMIN}. \\
\textbf{Gatilho} & O administrador acessa a tela administrativa de Gestão de Usuários. \\
\textbf{Fluxo Principal} & 
1. O sistema apresenta a tabela com exatamente dez usuários por página e busca.\\
& 2. O administrador aciona o botão "Novo Usuário" e preenche os dados: nome, e-mail institucional e perfil desejado (\texttt{ADMIN}, \texttt{PROFESSOR} ou \texttt{STUDENT}).\\
& 3. Se o perfil selecionado for \texttt{STUDENT}, o sistema exibe o campo obrigatório "Matrícula". Se for \texttt{ADMIN} ou \texttt{PROFESSOR}, o campo é ocultado.\\
& 4. O sistema gera uma senha provisória aleatória e define \texttt{must\_change\_password = true}.\\
& 5. O sistema persiste o registro na tabela \texttt{users} com hash BCrypt.\\
& 6. O sistema exibe as credenciais provisórias para fornecimento ao usuário. \\
\textbf{Fluxos Alternativos} & 
\textbf{[1a] Inativação/Reativação}: O administrador localiza um usuário e altera sua situação para inativo (\texttt{is\_active = false}), bloqueando logins futuros sem apagar dados históricos.\\
& \textbf{[1b] Edição de Dados}: O administrador atualiza o nome completo ou perfil. \\
\textbf{Fluxos de Exceção} & 
\textbf{[2e] Violação de Domínio}: Se o e-mail não terminar em \texttt{@academico.ufs.br}, o sistema rejeita a operação com erro de validação.\\
& \textbf{[3e] Ausência de Matrícula para Estudante}: O sistema impede o envio do formulário.\\
& \textbf{[5e] E-mail Duplicado}: O sistema detecta colisão de unicidade e impede o cadastro. \\
\textbf{Pós-condições} & Novo usuário criado e pronto para o primeiro acesso provisório. \\
\textbf{Regras de Negócio} & \textbf{RN01}, \textbf{RN03} (Matrícula Obrigatória Apenas para Estudante), \textbf{RN05}. \\
\textbf{Requisitos Associados} & \texttt{RF01}, \texttt{RF04}, \texttt{RF05}, \texttt{RNF01}. \\
\bottomrule
\end{tabularx}
\fonte{Elaborado pelo autor (2026).}
\end{table}

% ---------------------------------------------------------------------------
% UC05
% ---------------------------------------------------------------------------
\subsection{UC05: Parametrizar Módulos e Gatilhos Clínicos GTT}
\label{uc05_especificacao}

\begin{table}[!ht]
\centering
\caption{Especificação Detalhada do Caso de Uso UC05}
\label{tab_uc05}
\footnotesize
\begin{tabularx}{\textwidth}{p{3.5cm} X}
\toprule
\textbf{Propriedade} & \textbf{Especificação Técnica e Comportamental} \\
\midrule
\textbf{Identificador e Nome} & \texttt{UC05: Parametrizar Módulos e Gatilhos Clínicos GTT} \\
\textbf{Atores} & \texttt{ADMIN} (Ator Primário). \\
\textbf{Nível de Meta} & Nível Usuário (\textit{Sea Level}). \\
\textbf{Pré-condições} & Usuário autenticado como \texttt{ADMIN}. \\
\textbf{Gatilho} & O administrador acessa a tela de Parametrização do Catálogo IHI-GTT. \\
\textbf{Fluxo Principal} & 
1. O sistema exibe os seis módulos canônicos do IHI com seus códigos.\\
& 2. O administrador seleciona um módulo e visualiza os gatilhos vinculados.\\
& 3. O administrador cadastra ou atualiza um gatilho informando: código (ex.: \texttt{"M03"}), nome oficial, descrição clínica e status ativo/inativo.\\
& 4. O sistema valida a unicidade do código do gatilho no banco de dados.\\
& 5. O sistema persiste a alteração e atualiza o catálogo institucional. \\
\textbf{Fluxos de Exceção} & 
\textbf{[4e] Código de Gatilho Duplicado}: O sistema impede a duplicação.\\
& \textbf{[5e] Tentativa de Exclusão Física de Gatilho Vinculado}: A integridade referencial \texttt{ON DELETE RESTRICT} barra a exclusão caso o gatilho já esteja referenciado em auditorias anteriores, sugerindo a sua inativação lógica. \\
\textbf{Pós-condições} & Catálogo clínico IHI parametrizado e disponível para turmas e atividades. \\
\textbf{Regras de Negócio} & Integridade taxonômica conforme padrão IHI \cite{ihi2009}. \\
\textbf{Requisitos Associados} & \texttt{RF06}, \texttt{RF07}, \texttt{RNF04}. \\
\bottomrule
\end{tabularx}
\fonte{Elaborado pelo autor (2026).}
\end{table}

% ---------------------------------------------------------------------------
% UC06
% ---------------------------------------------------------------------------
\subsection{UC06: Consultar Catálogo Didático de Gatilhos e Gravidades}
\label{uc06_especificacao}

\begin{table}[!ht]
\centering
\caption{Especificação Detalhada do Caso de Uso UC06}
\label{tab_uc06}
\footnotesize
\begin{tabularx}{\textwidth}{p{3.5cm} X}
\toprule
\textbf{Propriedade} & \textbf{Especificação Técnica e Comportamental} \\
\midrule
\textbf{Identificador e Nome} & \texttt{UC06: Consultar Catálogo Didático de Gatilhos e Gravidades} \\
\textbf{Atores} & \texttt{STUDENT}, \texttt{PROFESSOR}, \texttt{ADMIN} (Atores Primários). \\
\textbf{Nível de Meta} & Nível Usuário (\textit{Sea Level}). \\
\textbf{Pré-condições} & Usuário autenticado na plataforma. \\
\textbf{Gatilho} & O usuário acessa a opção "Catálogo de Gatilhos" no menu lateral. \\
\textbf{Fluxo Principal} & 
1. O sistema renderiza os cinquenta e três gatilhos agrupados por módulo.\\
& 2. O usuário pode filtrar por módulo temático ou digitar na barra de busca.\\
& 3. O usuário clica em um gatilho para expandir o modal de detalhes clínicos.\\
& 4. O sistema apresenta a descrição técnica do gatilho, exemplos práticos de ocorrência em prontuários e a faixa típica de dano na escala NCC MERP.\\
& 5. O usuário consulta as nove categorias de gravidade (A até I) e os critérios de diferenciação entre dano temporário, permanente ou óbito. \\
\textbf{Fluxos Alternativos} & 
\textbf{[2a] Busca com Debounce}: À medida que o usuário digita, o sistema aplica filtro reativo no frontend sem sobrecarregar o backend com requisições redundantes. \\
\textbf{Pós-condições} & Apoio pedagógico e consulta didática realizada sem mutação de estado. \\
\textbf{Regras de Negócio} & Classificação NCC MERP \cite{nccmerp2001}. \\
\textbf{Requisitos Associados} & \texttt{RF08}, \texttt{RF09}, \texttt{RNF03}. \\
\bottomrule
\end{tabularx}
\fonte{Elaborado pelo autor (2026).}
\end{table}
